% =====================================================================================================
\section{Sample construction: sources and observations}\label{sec:samples}
% =====================================================================================================

Sample construction runs in two stages. First, \emph{sources} are chosen and labelled from published tables. Second, for each
source, \emph{observations} are chosen from the archive catalogue by a rule that uses only times, exposures, count rates and
positions, never the spectrum or the light curve. Data are downloaded and processed only after both lists are written. The
processing can reject an observation (quality rules); the next candidate of the same time group then takes its place.
\Cref{fig:selectflow} summarises the procedure.

\begin{figure}[htbp]
\centering
% Source labelling, eligibility and time-stratified sampling (Algorithm 1). Copied from report_en/figures/flowcharts (styles in workflow_report.tex).
\begin{tikzpicture}[x=1cm, y=1cm,
  st/.style={fc proc, text width=2.62cm, minimum height=15mm},
  lab/.style={font=\tiny, text=inkgrey, inner sep=1pt}]
  % row 1: sources and eligible pointings
  \node[fc data, text width=2.62cm, minimum height=15mm] (lb) at (1.45,0) {label lists of the version (Table~\ref{tab:labelrules})};
  \node[st] (mt) at (4.62,0) {positions matched to RXTE catalogues ($\le0.1^\circ$) or \texttt{nicermastr} ($\le0.05^\circ$)};
  \node[st] (el) at (7.79,0) {eligible pointings: exposure $\ge1$\,ks, rate, epoch; source needs a minimum number};
  \node[st] (gr) at (10.96,0) {sorted by time and split into $K$ groups (RXTE 15, NICER 8)};
  \node[st] (rk) at (14.13,0) {within a group: random order from a seeded generator};
  \foreach \a/\b in {lb/mt,mt/el,el/gr,gr/rk} \draw[fc arr] (\a) -- (\b);
  % row 2: the try-and-accept loop (right to left)
  \node[fc data, text width=2.62cm, minimum height=15mm] (dl) at (14.13,-2.5) {download and process the next candidate};
  \node[fc dec, text width=1.7cm, aspect=1.6] (q) at (10.6,-2.5) {quality rules pass?};
  \node[fc, draw=inkgrey, fill=boxgrey, text width=2.62cm, minimum height=15mm] (ok) at (6.95,-2.5) {accept; go to the next group};
  \node[fc dec, text width=1.55cm, aspect=1.6] (mx) at (10.6,-4.75) {tries left (${\le}M$)?};
  \node[fc, draw=inkgrey, fill=white, dashed, text width=2.62cm, minimum height=11mm] (gv) at (6.95,-4.75) {group left empty};
  \draw[fc arr] (rk) -- (dl);
  \draw[fc arr] (dl) -- (q);
  \draw[fc arr] (q) -- node[lab, above] {yes} (ok);
  \draw[fc arr] (q) -- node[lab, right] {no} (mx);
  \draw[fc arr] (mx.east) -| node[lab, pos=0.25, below] {yes} (dl.south);
  \draw[fc arr] (mx) -- node[lab, above] {no} (gv);
  \node[fc note, text width=5.2cm, anchor=north west] at (0.15,-1.85)
     {$M$ = 6 tries per group for RXTE, 3 for NICER; seeds \code{SEED + position} (v11: $+1000$, v13: $+2000$; Table~\ref{tab:seeds}). v4b1 used every eligible pointing instead.};
\end{tikzpicture}

\caption{Selection of sources and observations. Apart from the quality rules, the choice of an observation depends only on
the label list, the catalogue position, the eligibility rules, the time order and a seeded random order.}\label{fig:selectflow}
\end{figure}

% -----------------------------------------------------------------------------------------------------
\subsection{Source labels}\label{sec:labels}
The label is a property of the source: $y=1$ for a black hole and $y=0$ for a neutron star. A BH needs a dynamical mass
function, and an NS needs type-I X-ray bursts or coherent accretion-powered pulsations. BH candidates (X-ray evidence only)
are never labels. \Cref{tab:labelrules} lists the rule of each round. The realised source lists, with the evidence and a
reference URL for every source, are \file{data/<v>/sources.csv}.

\begin{table}[htbp]
\centering\small
\caption{Labelling and eligibility rules by round. ``Pointing'' is the largest separation allowed between an observation's
pointing and the source position. Every round also requires a catalogue exposure $\ge1000$\,s (RXTE: and
\code{STD1RATE}~$\ge5$\,c\,s$^{-1}$\,PCU$^{-1}$).}\label{tab:labelrules}
\begin{tabularx}{\textwidth}{@{}>{\raggedright\arraybackslash}p{1.35cm}>{\raggedright\arraybackslash}X>{\raggedright\arraybackslash}X>{\raggedright\arraybackslash}p{3.9cm}@{}}
\toprule
Round & BH rule & NS rule & Pointing; time window; minimum eligible \\
\midrule
v1 & 4 BlackCAT dynamical BHs & 4 Galloway+2008 bursters & $0.1^\circ$ from the catalogue median; gain epoch 5 to MJD 54094; 30 \\
v2, v3 & every BlackCAT Table~4 BH with a MissionLongData catalogue & every Galloway+2008 appendix-A burster with a catalogue & $0.1^\circ$; gain epoch 5 (MJD 51677--55931); 10 \\
v3c & 16 BlackCAT candidates (prediction targets only) & -- & as v2 \\
v4b2 & BlackCAT BHs found by position & Galloway+2008 Table~3 bursters found by position & $0.1^\circ$ from Sesame; all gain epochs; 10 \\
v6 & HMXB BHs (Cyg X-1, LMC X-3) & accreting millisecond pulsars; 4 slow pulsars as a stress set & as v4b2 \\
v7b2 & Cyg X-1, LMC X-1, LMC X-3 \citep{Marcel2026} & -- & $0.1^\circ$ from Sesame; gain epoch 5; 10 \\
v7c (MAXI) & \citet{Marcel2026} Table~1 & MINBAR bursters without pulsations; pulsars as a separate class & MAXI source within $0.1^\circ$ ($0.3^\circ$ fallback); $\ge100$ daily points \\
v8b & Marcel+2026 BHs not used before (GS 1354$-$64, SS 433) & -- & $0.1^\circ$ from Sesame; all gain epochs; 3 \\
v9--v13 (NICER) & Marcel+2026 Table~1, Sesame positions & MINBAR bursters; entries with another MINBAR entry within $0.05^\circ$ removed & target within $0.05^\circ$; mission; 5 \\
\bottomrule
\end{tabularx}
\end{table}

\paragraph{Confused fields.} NICER does not image, so a MINBAR source with another MINBAR source within $0.05^\circ$ is
flagged \code{excluded_confused} and not used. Nine of the 115 MINBAR sources are removed this way in \file{47_v9_select.py}.
For RXTE, the $0.1^\circ$ pointing rule plays the same role; \file{config.py} notes, for example, that 4U 1907+097 lies
about $0.4^\circ$ from the candidate XTE J1908+094.

% -----------------------------------------------------------------------------------------------------
\subsection{Catalogue-level eligibility}\label{sec:eligibility}
For RXTE (\file{03_select_observations.py} and its copies), each source's MissionLongData catalogue goes through the steps
below. The number of rows left after each step is written to \file{results/<v>/attrition_catalog.csv}:
\begin{enumerate}
\item drop duplicated \code{OBSID}s;
\item keep the time window of the round: gain epoch 5 (MJD 51677 to 54094 in v1, to 55931 from v2), or all gain epochs
(v4b2, v6, v8b);
\item keep \code{EXPOSURE}~$\ge1000$\,s;
\item keep \code{STD1RATE}~$\ge5$\,c\,s$^{-1}$\,PCU$^{-1}$;
\item keep pointings within $0.1^\circ$ of the source position (great-circle distance);
\item drop the source if fewer than the minimum number of pointings remain (\cref{tab:labelrules}). The source stays in
\file{sources.csv} with \code{included = False}.
\end{enumerate}
For NICER (\file{47_v9_select.py}), the TAP query already restricts the target position to $0.05^\circ$. The remaining
steps are: drop rows without a time, drop duplicated ObsIDs, keep \code{exposure}~$\ge1000$\,s, and require at least 5
eligible observations and a non-confused field.

% -----------------------------------------------------------------------------------------------------
\subsection{Time-stratified ranking}\label{sec:ranking}
The eligible observations of a source are sorted by time and split into $K$ consecutive groups of nearly equal size
(\code{numpy.array_split}): $K=15$ for RXTE, 8 for NICER. A seeded generator puts the observations of each group in a
random order (\cref{alg:sampling}). The seed is \code{SEED} plus an offset plus the 1-based position of the source in the
source list (\cref{tab:seeds}). The result is \file{data/<v>/observation_candidates.csv}, with one row per eligible
observation, its group (\code{time_bin}) and its rank in the group (\code{rank_in_bin}). This file is written before any
data file is downloaded.

The stratification spreads each source's sample over its history, across outbursts and states. Because the order is
random within a group, the choice does not depend on the data. If a source has fewer eligible observations than groups,
some groups are empty and every observation is a candidate.

\begin{algorithm}[htbp]
\caption{Time-stratified ranking and acceptance. RXTE: \file{03_select_observations.py} with
\file{04_fetch_process.py}; NICER: \file{47}/\file{55}/\file{60} with \file{48}/\file{52}/\file{56}/\file{61}.}\label{alg:sampling}
\begin{algorithmic}[1]
\Require eligible observations $\{(o_i,t_i)\}$ of one source; number of groups $K$; seed $\sigma$; tries per group $M$
\State sort by $t_i$ and split the index list into $K$ consecutive groups $G_1,\dots,G_K$
\State $\mathrm{rng}\gets\mathrm{default\_rng}(\sigma)$
\For{$b=1,\dots,K$}
  \State $\pi\gets\mathrm{rng.permutation}(G_b)$ \Comment{rank within the group; written to the candidate table}
\EndFor
\Statex \textit{--- later, in the fetch script, for each group independently ---}
\For{$r=1,\dots,\min(M,|G_b|)$}
  \State download and process $o_{\pi_r}$; record status and reason
  \If{every quality rule passes} accept $o_{\pi_r}$; \textbf{break} \EndIf
\EndFor
\end{algorithmic}
\end{algorithm}

\begin{table}[htbp]
\centering\small
\caption{Sampling parameters. ``Position'' is the 1-based position of the source in the list processed by the script
(for NICER, \file{data/v9/sources.csv}, counting excluded sources too). \code{SEED = 42}.}\label{tab:seeds}
\begin{tabular}{@{}l l l l@{}}
\toprule
Sample & Groups $K$ & Seed $\sigma$ & Tries per group $M$ \\
\midrule
RXTE v1, v2, v3c, v7b2, v8b & 15 & \code{SEED + position} & 6 \\
RXTE v4b2 (new sources) & 15 & \code{SEED + 100 + i} & 6 \\
RXTE v4b1 & all eligible & -- & 1 \\
NICER v9 (N1) & 8 & \code{SEED + position} & 3 \\
NICER v11 (N2) & 8 & \code{SEED + 1000 + position} & 3 \\
NICER v13 (N3) & 8 & \code{SEED + 2000 + position} & 3 \\
\bottomrule
\end{tabular}
\end{table}

% -----------------------------------------------------------------------------------------------------
\subsection{Acceptance loop and quality rules}\label{sec:quality}
The fetch scripts work through the (source, group) pairs in parallel: 4 threads for RXTE; for NICER, 8 threads in v9 and
2 in v11 and v13. Within a pair the candidates are tried in rank order. Every tried candidate becomes one row of
\file{data/<v>/observations.csv}, with its status, its reason, the download URL and the quality quantities. Rejected
observations are therefore documented, not lost.

\begin{table}[htbp]
\centering\small
\caption{Post-download quality rules. An observation is accepted only if no rule fires. ``Reason'' is the string written
to \file{observations.csv}.}\label{tab:qualityrules}
\begin{tabularx}{\textwidth}{@{}l l >{\raggedright\arraybackslash}X@{}}
\toprule
Data & Reason & Rule \\
\midrule
RXTE PCA & \code{fetch_or_read_error} & a file is missing in every AO directory, or cannot be read \\
 & \code{nonfinite} & a rebinned rate or error is not finite \\
 & \code{short_exposure} & Standard-2 \code{EXPOSURE}~$<1000$\,s \\
 & \code{low_snr} & net 5--25\,keV rate / its error $<10$ (\cref{eq:F}) \\
 & \code{high_deadtime} & dead-time fraction $>0.10$ (\cref{eq:dtf}) \\
 & \code{coord_mismatch} & header pointing $>0.1^\circ$ from the source position \\
 & \code{nonstandard_channels} & spectrum or response does not have 129 channels \\
 & \code{ebounds_coverage} & the response's energy range does not cover the common grid \\
\midrule
NICER XTI & \code{error: ...} & download or parsing error (after retries) \\
 & \code{missing (event file not found)} & no file in the catalogue month or the months before and after \\
 & \code{missing (< 3 valid segments)} & fewer than 3 usable 128\,s segments (\cref{sec:nicer}) \\
 & \code{excluded (2-10 keV rate x < 20.0)} & mean 2--10\,keV rate in the valid segments $<20$\,c\,s$^{-1}$ \\
\bottomrule
\end{tabularx}
\end{table}

Two further screens act after acceptance and do not trigger a replacement. Type-I bursts are flagged, or the affected
segments removed (\cref{sec:bursts}). From v12, NICER observations are screened on their 3C50 background fraction
(\cref{sec:3c50}).

% -----------------------------------------------------------------------------------------------------
\subsection{Disjoint NICER observation sets}\label{sec:disjoint}
NICER was used three times, each time with observations that had never been tried before:
\begin{itemize}
\item \textbf{N1} (v9, re-read in full in v10): the first time-stratified sample of the v9 sources.
\item \textbf{N2} (v11): for each included v9 source, the eligible observations \emph{not} among the 792 v9 tries, ranked
with \code{SEED + 1000 + position}. Sources not included in v9 but with 1--4 eligible observations and a non-confused field
contribute every eligible observation, one per group.
\item \textbf{N3} (v13): eligible observations of the included v9 sources that are among neither the 792 v9 tries nor the
588 v11 tries, ranked with \code{SEED + 2000 + position}.
\end{itemize}
The selection scripts compute the overlap with the earlier tries and stop with exit code 3 if it is not empty (check IC3 in
v11). All three sets come from the same cached \code{nicermastr} tables, so no new catalogue query can change the
candidates.

% -----------------------------------------------------------------------------------------------------
\subsection{Resulting samples}\label{sec:samplesizes}
\Cref{tab:samples} lists the samples that enter models or tests. Counts are recomputed from the committed
\file{observations.csv} and processed tables.

{\small\setlength{\tabcolsep}{3.5pt}
\begin{longtable}{@{}>{\raggedright\arraybackslash}p{2.9cm} >{\raggedright\arraybackslash}p{2.8cm} r r >{\raggedright\arraybackslash}p{2.1cm} >{\raggedright\arraybackslash}p{4.3cm}@{}}
\caption{Samples. ``Tried'' counts every downloaded candidate; ``Obs.'' counts the accepted (or kept) observations;
``Sources'' gives the total and the BH/NS split.}\label{tab:samples}\\
\toprule
Sample & Round, instrument & Tried & Obs. & Sources & Use \\
\midrule\endfirsthead
\toprule Sample & Round, instrument & Tried & Obs. & Sources & Use \\ \midrule\endhead
\bottomrule\endfoot
pilot & v1, PCA & 120 & 120 & 8 (4/4) & pipeline test \\
S1 (development) & v2, PCA, epoch 5 & 462 & 456 & 31 (7/24) & training and LOSO in v2--v8, v10 \\
BH candidates & v3c, PCA & 210 & 210 & 14 (14/0) & prediction only \\
S3 (all pointings) & v4b1, PCA, epoch 5 & 7\,975 & 7\,949 & 31 (7/24) & observation-budget checks \\
v4b2 sources & v4b2, PCA, all epochs & 221 & 221 & 15 (2/13) & S2 = S1 $\cup$ v4b2 (46 sources, 677 obs.) \\
E\textsubscript{new} + stress & v6, PCA, all epochs & 180 & 180 & 12 (2/10) & 8 new sources (120 obs.); 4 slow pulsars (60 obs.) \\
E (external) & v6 & -- & 341 & 23 (4/19) & frozen-model test \\
persistent BHs & v7b2, PCA, epoch 5 & 45 & 45 & 3 (3/0) & Cyg X-1, LMC X-1, LMC X-3 \\
new BHs & v8b, PCA, all epochs & 26 & 23 & 2 (2/0) & GS 1354$-$64, SS 433 \\
X (external) & v8b & -- & 379 & 26 (7/19) & frozen-model test; 99 hard-like obs.\ (5/15 sources) \\
MAXI & v7c2, GSC & -- & 132\,057 & 100: 13 BH, 48 NPNS, 39 PSR & daily points (not observations); point-level classifiers \\
N1 & v9 (prefix), XTI & 792 & 368 & 63 (9/54) & first NICER sample \\
N1 full files & v10, XTI & 368 & 366 & 63 (9/54) & N1 re-read to 1\,GB \\
N2 & v11, XTI & 588 & 318 & 60 (9/51) & unseen observations \\
N1+N2, 3C50 & v12, XTI & -- & 534 & 58 (9/49) & corrected and screened; frozen training set for v13 \\
N3 & v13, XTI & 412 & 238 & 37 (8/29) & unseen observations \\
N3, 3C50 & v13, XTI & -- & 198 & 35 (8/27) & corrected and screened test set \\
\end{longtable}}
